\documentclass[11pt]{article}

\usepackage[margin=1in]{geometry}
\usepackage{amsmath, amssymb, amsthm, amsfonts}
\usepackage{mathtools}
\usepackage{bbm}
\usepackage{bm}
\usepackage{graphicx}
\usepackage{enumitem}
\usepackage{booktabs}
\usepackage{hyperref}
\usepackage{cleveref}
\usepackage{listings}
\usepackage{courier}

\lstdefinestyle{python}{
  language=Python,
  basicstyle=\footnotesize\ttfamily,
  keywordstyle=\bfseries,
  commentstyle=\itshape,
  showstringspaces=false,
  frame=single,
  breaklines=true
}

\title{%
Non--Euclidean Logic, Contractive Multiplicity Dynamics,\\
and Heavy--Tailed Forcing:\\
A Joint NEL--$\Lambda_m$--Cauchy Framework}

\author{%
\Large \textit{In legacy of Nicole Thorp} \\ \\ \Large Citizen Gardens \\ \& \\ \textit{The Prime Materia Commons}
}

\date{May 8, 2026}

\theoremstyle{plain}
\newtheorem{theorem}{Theorem}[section]
\newtheorem{proposition}[theorem]{Proposition}
\newtheorem{lemma}[theorem]{Lemma}
\newtheorem{corollary}[theorem]{Corollary}

\theoremstyle{definition}
\newtheorem{definition}[theorem]{Definition}
\newtheorem{assumption}[theorem]{Assumption}
\newtheorem{remark}[theorem]{Remark}

\numberwithin{equation}{section}

\begin{document}

\maketitle

\begin{abstract}
This report formalizes a joint framework combining Nicole Thorp's
Non--Euclidean Logic (NEL) with contractive multiplicity dynamics
$\Lambda_m$ and heavy--tailed (Cauchy / $\alpha$--stable) forcing.
The logical layer is modeled as a positively curved Riemannian manifold
$L$ of dimension four, endowed with a dynamic metric $M$ driven by a
stress functional. The dynamical layer is a contractive stochastic
recursion $x_{t+1} = T_t x_t + \xi_t$ on a finite--dimensional complex
state space, where $\rho(T_t) < 1$ and $\xi_t$ is either Gaussian or
Cauchy--type noise. Stress is defined via a negative--definite quadratic
resonance $R(x)$ and mapped into the NEL metric. We give a coherent
mathematical specification, a baseline convergence statement, and a
simulation scheme for comparing Gaussian and Cauchy forcing in the NEL
geometry. The document is written as a technical foundation for further
analysis and numerical validation.
\end{abstract}
\newpage
\tableofcontents
\newpage
\section{Executive summary}

This document specifies a coupled system with three principal
components:

\begin{enumerate}[label=(\roman*)]
  \item A Non--Euclidean logical manifold $L$ with positive curvature,
  dynamic metric $M(\cdot)$, and stress--dependent geodesic structure,
  following the NEL architecture proposed by Thorp.

  \item A $\Lambda_m$--tuned contractive recursion $x_{t+1} = T_t x_t + \xi_t$
  on $\mathbb{C}^n$, with negative--definite quadratic resonance
  $R(x) = - x^\ast Q x$ and Lyapunov candidate $V(x) = \Phi(-R(x))$.

  \item A noise model $\xi_t$ which is either Gaussian (finite variance)
  or Cauchy / $\alpha$--stable with $\alpha = 1$, inducing heavy tails
  in stress peaks and return times.
\end{enumerate}

The main design objective is to realize a system in which:

\begin{itemize}
  \item The logical geometry is given by NEL (positive curvature,
  dynamic metric, local and global warps).

  \item The state dynamics are globally stabilizing by spectral and
  Lyapunov constraints ($\Lambda_m$ / ROC style).

  \item The innovation process is heavy--tailed, so the statistics of
  extreme events and return maps are Cauchy--like rather than Gaussian,
  even though the state remains confined to a bounded region.
\end{itemize}

We give a minimal convergence result for a contractive recursion with
heavy--tailed forcing, a mapping from resonance to NEL stress, and a
simulation scheme for Gaussian--vs--Cauchy comparison.

\section{Non--Euclidean logical manifold}

\subsection{Logical manifold and coordinates}

\begin{definition}[Logical manifold]
Let $L$ be a $4$--dimensional smooth Riemannian manifold, with
underlying differentiable structure diffeomorphic to $S^4$ and metric
tensor $g_\sigma$ parametrized by a scalar stress variable
$\sigma \in [0,\infty)$. The pair $(L, g_\sigma)$ is called the
Non--Euclidean Logical Manifold.
\end{definition}

Following Thorp, each proposition $P$ is represented as a point
$x_P \in L$ with four distinguished coordinates
\[
x_P = (T, I, C, K) \in \mathbb{R}^4,
\]
interpreted as temporal relevance $T$, influence $I$, certainty $C$,
and contextual coherence $K$, embedded into $L$ via a smooth chart.
The curvature of $L$ is assumed strictly positive with respect to
$g_\sigma$ for all admissible $\sigma$.

\subsection{Dynamic metric}

\begin{definition}[Dynamic metric function]
The Dynamic Metric function
\[
M : L \times [0,\infty) \to \mathrm{Sym}^+_4
\]
assigns to each $(x,\sigma)$ a positive--definite matrix
$M(x,\sigma)$ such that $g_\sigma$ is locally represented as
\[
g_\sigma(x)(u,v) = u^\top M(x,\sigma)\, v,
\quad u,v \in \mathbb{R}^4.
\]
The dependence on $\sigma$ is such that increasing $\sigma$ increases
local curvature and decreases effective distances between selected
points.
\end{definition}

For minimal modeling, one may adopt a scalar metric factor
$m(\sigma) \ge 1$ and set
\[
M(x,\sigma) = m(\sigma)\, M_0(x),
\]
where $M_0$ is a baseline metric and $m$ is a scalar dynamic
multiplier.

\begin{assumption}[Monotonic curvature scaling]
The function $m : [0,\infty) \to [1,\infty)$ is continuous and
nondecreasing, with $m(0) = 1$.
\end{assumption}

This captures the intuition that higher stress collapses distances,
increasing curvature.

\subsection{Stress--induced warping}

Let $d_0(\cdot,\cdot)$ denote the Riemannian distance induced by $M_0$ and
$d_\sigma(\cdot,\cdot)$ the distance induced by $M(\cdot,\sigma)$.
Under the scalar metric factor assumption,
\begin{equation}
  d_\sigma(x,y) = \frac{1}{\sqrt{m(\sigma)}}\, d_0(x,y).
  \label{eq:distance-warp}
\end{equation}

Define two thresholds
\[
0 < \sigma_{\text{local}} < \sigma_{\text{global}}.
\]
For $\sigma < \sigma_{\text{local}}$ only mild ``local'' curvature
adjustments occur; for $\sigma_{\text{local}} \le \sigma <
\sigma_{\text{global}}$ one obtains localized dimpling; for
$\sigma \ge \sigma_{\text{global}}$ the manifold undergoes a near--global
warp.

\section{Contractive multiplicity dynamics}

\subsection{State space and update law}

Let $n \in \mathbb{N}$ and consider a complex state
$x_t \in \mathbb{C}^n$ evolving in discrete time.

\begin{definition}[Λ$_m$--tuned ROC recursion]
A $\Lambda_m$--tuned ROC recursion is a stochastic process
$(x_t)_{t\ge0}$ satisfying
\begin{equation}
  x_{t+1} = T_t x_t + \xi_t,
  \label{eq:roc-recursion}
\end{equation}
where
\begin{enumerate}[label=(\alph*)]
  \item $T_t \in \mathbb{C}^{n\times n}$ is a random or
  time--varying linear operator with
  \[
  \rho(T_t) \le 1 - \eta
  \]
  for some fixed $\eta \in (0,1)$ and all $t$;

  \item $\xi_t \in \mathbb{C}^n$ is an innovation noise sequence
  (Gaussian or Cauchy--type), independent of $x_0$;

  \item The pair $(T_t,\xi_t)$ is adapted to a filtration
  $(\mathcal{F}_t)$ and satisfies standard measurability constraints.
\end{enumerate}
\end{definition}

The spectral constraint is the Λ$_m$ contractivity condition.

\subsection{Quadratic resonance and Lyapunov functional}

\begin{definition}[Resonance and Lyapunov candidate]
Let $Q \in \mathbb{C}^{n\times n}$ be Hermitian positive--definite.
Define the resonance
\begin{equation}
  R(x) = - x^\ast Q x \le 0
\end{equation}
and the Lyapunov candidate
\begin{equation}
  V(x) = \Phi(-R(x)) = \Phi(x^\ast Q x),
\end{equation}
where $\Phi : [0,\infty) \to [0,\infty)$ is strictly increasing and
$C^1$.
\end{definition}

A natural choice is $\Phi(s) = s$ or $\Phi(s) = \log(1+s)$.

\begin{assumption}[Drift condition]
There exists $\delta > 0$ and $b \ge 0$ such that
\begin{equation}
  \mathbb{E}[V(x_{t+1}) \mid x_t] \le (1-\delta)\, V(x_t) + b
  \label{eq:lyap-drift}
\end{equation}
whenever $\|x_t\|$ is sufficiently large.
\end{assumption}

This is a Foster--Lyapunov condition ensuring tightness of $(x_t)$
under suitable regularity conditions.

\subsection{Induced stress variable}

To couple the ROC dynamics with NEL, define a scalar stress
$\sigma_t$ as a function of $x_t$.

\begin{definition}[Stress mapping]
Let $h : [0,\infty) \to [0,\infty)$ be continuous and increasing.
Define
\begin{equation}
  \sigma_t = h\bigl(x_t^\ast Q x_t\bigr).
  \label{eq:stress-from-roc}
\end{equation}
\end{definition}

One may set $h(s) = s$ or $h(s) = s / (1+s)$ depending on scaling
requirements.

\section{Noise models: Gaussian vs Cauchy}

\subsection{Gaussian innovations}

\begin{definition}[Gaussian noise]
The sequence $(\xi_t)$ is Gaussian if
$\xi_t \sim \mathcal{N}(0,\Sigma)$ in $\mathbb{R}^{2n}$ (identified
with $\mathbb{C}^n$), with finite covariance $\Sigma$ and
$\xi_t$ i.i.d.
\end{definition}

In this case, all components have finite moments of all orders and
standard central limit behavior is expected for linear functionals.

\subsection{Cauchy / $\alpha$--stable innovations}

\begin{definition}[Symmetric Cauchy noise]
The sequence $(\xi_t)$ is symmetric Cauchy if each component
$\xi^{(k)}_t$ has density
\begin{equation}
  f(x) = \frac{1}{\pi}\, \frac{\gamma}{x^2 + \gamma^2},
\end{equation}
for some scale parameter $\gamma > 0$, with components independent
across $k$ and $t$.
\end{definition}

\begin{remark}
Cauchy distributions are $1$--stable and heavy--tailed: they admit no
finite mean or variance, and sums scaled by $n$ converge in law to a
Cauchy rather than a Gaussian limit.%
\footnote{For standard properties of the Cauchy distribution and
heavy--tailed stable laws, see e.g.\ \cite{cauchy-notes,heavy-tails}.}
\end{remark}

In this regime, the Lyapunov drift condition
\eqref{eq:lyap-drift} must control the state despite unbounded second
moments of the innovations.

\section{Coupled NEL--$\Lambda_m$ system}

\subsection{Coupling scheme}

The coupled system is defined by the following components:

\begin{itemize}
  \item ROC recursion \eqref{eq:roc-recursion} on $x_t \in \mathbb{C}^n$.

  \item Stress variable $\sigma_t$ defined by
  \eqref{eq:stress-from-roc}.

  \item NEL manifold $(L,g_{\sigma_t})$ with metric factor
  $m(\sigma_t)$ and distance warping
  \eqref{eq:distance-warp}.
\end{itemize}

The NEL manifold influences the ROC dynamics through the choice of
$T_t$ and possibly through the distribution of $\xi_t$,
if one allows $T_t$ and $\xi_t$ to depend on geometric quantities
(e.g.\ proximity to fixed points or Crisis Core).

\subsection{Minimal coupled model}

A minimal coupled model is obtained by:
\begin{enumerate}[label=(\alph*)]
  \item Choosing a constant contraction $T \in \mathbb{R}^{n\times n}$
  with eigenvalues in $(0,1)$.

  \item Fixing $Q \succ 0$ and $h(s) = s$ so that
  $\sigma_t = x_t^\ast Q x_t$.

  \item Setting $m(\sigma) = 1 + k\sigma$ for some $k > 0$.

  \item Allowing $\xi_t$ to be either Gaussian or Cauchy.
\end{enumerate}

In this simplified case the geometry does not feed back into
$T$ directly, but the stress variable and metric factor can be logged
and analyzed as functions of time.

\subsection{Convergence and heavy tails}

The following theorem summarizes the qualitative behavior of the
minimal model.

\begin{theorem}[Qualitative behavior under Gaussian vs Cauchy forcing]
\label{thm:gauss-cauchy-qualitative}
Consider the minimal coupled model with fixed $T$, fixed $Q \succ 0$
and $m(\sigma) = 1 + k\sigma$.

\begin{enumerate}[label=(\roman*)]
  \item Suppose $(\xi_t)$ are i.i.d.\ Gaussian with finite covariance.
  Then, under standard minorization and irreducibility conditions, the
  Markov chain $(x_t)$ admits a unique stationary distribution with
  finite moments of all orders; linear functionals of $x_t$ satisfy a
  central limit theorem and the stationary distribution of
  $\sigma_t = x_t^\ast Q x_t$ has exponentially decaying tails.

  \item Suppose instead that $(\xi_t)$ are i.i.d.\ symmetric Cauchy
  with nondegenerate scale. Then $(x_t)$ admits at least one stationary
  distribution, but this stationary law is heavy--tailed: for generic
  nonzero linear functionals $u^\ast x_t$ the marginal distributions
  have power--law tails with asymptotic decay of order
  $|z|^{-1-\alpha}$ with $\alpha \approx 1$. In particular,
  $\mathbb{E}[\|x_t\|^2] = \infty$ and the stationary distribution of
  $\sigma_t$ is heavy--tailed.
\end{enumerate}
\end{theorem}

\begin{remark}
A rigorous proof requires standard Markov chain stability theory on
$\mathbb{R}^{2n}$ together with properties of affine recursions with
stable noise. The result is stated here as a design guideline rather
than a fully proved theorem. In both regimes, the contraction
$\rho(T)<1$ prevents almost sure divergence, but in the Cauchy case
the invariant law has no finite second moment.
\end{remark}

\section{Simulation scheme}

\subsection{Discrete approximation}

For numerical illustration, one may restrict to $n=3$ and choose
\[
T = \mathrm{diag}(a_1,a_2,a_3),
\quad 0 < a_i < 1,
\]
\[
Q = I_3,
\]
and implement the recursion
\begin{equation}
x_{t+1} = T x_t + \xi_t,
\quad
\sigma_t = \|x_t\|^2,
\quad
m_t = 1 + k \sigma_t.
\end{equation}

Two cases are run:
\begin{enumerate}[label=(\alph*)]
  \item Gaussian: $\xi_t \sim \mathcal{N}(0,\sigma^2 I_3)$.
  \item Cauchy: each component of $\xi_t$ is symmetric Cauchy with
    scale $\gamma$.
\end{enumerate}

For each case, one records:
\begin{itemize}
  \item time series of $\sigma_t$ and $m_t$,
  \item empirical tail distribution of $\sigma_t$,
  \item distribution of recovery times to a neighborhood of the origin
    after threshold crossings $\sigma_t > \theta$.
\end{itemize}

\subsection{Reference Python implementation}

The following Python snippets realize the minimal scheme.

\subsubsection*{Core dynamics}

\begin{lstlisting}[style=python,caption={ROC dynamics and NEL metric}]
import numpy as np

class ROCDynamics:
    def __init__(self, T, Q):
        self.T = T      # contraction matrix (n x n)
        self.Q = Q      # positive-definite matrix (n x n)

    def step(self, x, xi):
        return self.T @ x + xi

    def resonance(self, x):
        # R(x) = - x^* Q x
        return -np.conjugate(x).T @ self.Q @ x

class NELManifold:
    def __init__(self, k=1.0):
        self.k = k      # curvature scaling parameter

    def metric_factor(self, stress):
        # m(stress) = 1 + k * stress
        return 1.0 + self.k * stress
\end{lstlisting}

\subsubsection*{Noise models}

\begin{lstlisting}[style=python,caption={Gaussian and Cauchy noise}]
def gaussian_noise(dim, sigma):
    return np.random.normal(loc=0.0, scale=sigma, size=dim)

def cauchy_noise(dim, gamma):
    u = np.random.uniform(low=0.0, high=1.0, size=dim)
    return gamma * np.tan(np.pi * (u - 0.5))
\end{lstlisting}

\subsubsection*{Simulation loop}

\begin{lstlisting}[style=python,caption={Gaussian vs Cauchy simulation}]
def run_simulation(T, Q, steps, noise_type="gaussian",
                   sigma=0.1, gamma=0.1, k=1.0):
    n = T.shape[0]
    roc = ROCDynamics(T=T, Q=Q)
    nel = NELManifold(k=k)

    x = np.zeros(n, dtype=np.float64)
    stress_hist = []
    metric_hist = []

    for _ in range(steps):
        if noise_type == "gaussian":
            xi = gaussian_noise(n, sigma)
        elif noise_type == "cauchy":
            xi = cauchy_noise(n, gamma)
        else:
            raise ValueError("unknown noise_type")

        x = roc.step(x, xi)
        R = roc.resonance(x).real
        stress = max(0.0, -R)
        m = nel.metric_factor(stress)

        stress_hist.append(stress)
        metric_hist.append(m)

    return np.array(stress_hist), np.array(metric_hist)
\end{lstlisting}

This code can be used to generate empirical evidence for the claims of
\Cref{thm:gauss-cauchy-qualitative} by comparing tail behavior of
$\sigma_t$ under the two noise regimes.

\section{Extensions and open directions}

\subsection{Prime--weighted aggregation}

An orthogonal axis, not developed here, is to incorporate
prime--weighted fuzzy aggregation operators for combining multiple
ROC fibres into a single stress observable. This can be realized by
introducing a multiplicity--aware operator
\[
F(x_1,\dots,x_m; \omega_1,\dots,\omega_m),
\]
where $\omega_i$ are prime--based weights and
$x_i$ are fibre--level stress measures, with $F$ bounded in $[0,1]$
and satisfying t--norm properties. The output then drives $\sigma_t$.

\subsection{Feedback from NEL geometry}

More advanced coupled models can allow $T_t$ and even the distribution
of $\xi_t$ to depend on geometric quantities:
\[
T_t = T(g_{\sigma_t}, x_t), \quad
\xi_t \sim \mathcal{L}(g_{\sigma_t}, x_t),
\]
while preserving spectral and Lyapunov constraints. This leads to
nonautonomous dynamical systems on manifolds and requires tools from
Lyapunov stability on Riemannian manifolds.

\section{Conclusion}

We have specified a coupled framework in which:
\begin{itemize}
  \item NEL provides a positively curved logical manifold with
  stress--dependent metric;
  \item $\Lambda_m$ / ROC dynamics ensure contractive evolution of a
  finite--dimensional state via a negative--definite resonance and
  Lyapunov drift;
  \item Cauchy / $\alpha$--stable innovations yield heavy--tailed
  stress statistics and return maps, contrasting with Gaussian
  behavior.
\end{itemize}
This structure is intended as a rigorous foundation for further
analysis and numerical experimentation in the spirit of Nicole Thorp's
Non--Euclidean Logic program.

\appendix
\section*{Mathematical Appendix}

\section{Operator norm and spectral radius bounds}

\subsection{Basic inequalities}

Let $T \in \mathbb{C}^{n\times n}$ and denote by
$\|\cdot\|_2$ the spectral (operator) norm induced by the Euclidean
norm on $\mathbb{C}^n$. Denote by $\rho(T)$ the spectral radius.

\begin{lemma}[Spectral radius bound]
\label{lem:spectral-radius-bound}
For any $T \in \mathbb{C}^{n\times n}$,
\begin{equation}
  \rho(T) \le \|T\|_2.
\end{equation}
\end{lemma}

\begin{proof}
Let $\lambda$ be an eigenvalue of $T$ with eigenvector $v \ne 0$,
so $T v = \lambda v$. Then
\[
\|T v\|_2 = \|\lambda v\|_2 = |\lambda| \,\|v\|_2.
\]
On the other hand,
\[
\|T v\|_2 \le \|T\|_2 \,\|v\|_2.
\]
Combining these two relations gives $|\lambda| \le \|T\|_2$ for every
eigenvalue $\lambda$. Taking the supremum over all eigenvalues yields
$\rho(T) \le \|T\|_2$.
\end{proof}

\begin{corollary}[Uniform contraction from norm bound]
\label{cor:uniform-contraction}
Assume there exists $\kappa \in (0,1)$ such that
$\|T_t\|_2 \le \kappa$ for all $t \ge 0$. Then for any $x \in \mathbb{C}^n$,
\begin{equation}
  \|T_t x\|_2 \le \kappa \,\|x\|_2,
  \quad \forall t \ge 0.
\end{equation}
Consequently, for any $s < t$,
\begin{equation}
  \|T_{t-1} \cdots T_s x\|_2 \le \kappa^{t-s}\,\|x\|_2.
\end{equation}
\end{corollary}

\begin{proof}
The first statement is immediate from the definition of the operator
norm. The second follows by repeated application:
\[
\|T_{t-1}\cdots T_s x\|_2
  \le \|T_{t-1}\|_2 \cdots \|T_s\|_2 \,\|x\|_2
  \le \kappa^{t-s}\,\|x\|_2.
\]
\end{proof}

\subsection{Contractive ROC dynamics}

Consider the recursion
\begin{equation}
  x_{t+1} = T_t x_t + \xi_t,
  \label{eq:app-recursion}
\end{equation}
where $(T_t)$ and $(\xi_t)$ satisfy the conditions of the main text:
\begin{equation}
  \|T_t\|_2 \le \kappa < 1,
  \quad t \ge 0,
  \label{eq:app-contractive}
\end{equation}
and $(\xi_t)$ is a sequence of i.i.d.\ random vectors in
$\mathbb{C}^n$ (Gaussian or Cauchy / $1$--stable).

\begin{proposition}[Explicit bound on trajectories]
\label{prop:trajectory-bound}
Let $x_0 \in \mathbb{C}^n$ be arbitrary and assume
\eqref{eq:app-recursive-bound}--\eqref{eq:app-contractive}. Then for
each $t \ge 0$,
\begin{equation}
  \|x_t\|_2
  \le
  \kappa^t \|x_0\|_2 + 
  \sum_{j=0}^{t-1} \kappa^{t-1-j} \|\xi_j\|_2.
  \label{eq:xt-operator-bound}
\end{equation}
\end{proposition}

\begin{proof}
We iterate \eqref{eq:app-recursion}:
\begin{equation*}
  x_t
  = T_{t-1} \cdots T_0 x_0 + 
    \sum_{j=0}^{t-1} T_{t-1} \cdots T_{j+1} \xi_j.
\end{equation*}
Taking norms and using \Cref{cor:uniform-contraction},
\begin{align*}
  \|x_t\|_2
  &\le \|T_{t-1} \cdots T_0 x_0\|_2
    + \sum_{j=0}^{t-1} \|T_{t-1}\cdots T_{j+1}\xi_j\|_2 \\
  &\le \kappa^t \|x_0\|_2
    + \sum_{j=0}^{t-1} \kappa^{t-1-j} \|\xi_j\|_2,
\end{align*}
which is \eqref{eq:xt-operator-bound}.
\end{proof}

This decomposition is used repeatedly in the subsequent moment and
tail analyses.

\section{Quadratic resonance and Lyapunov drift}

\subsection{Quadratic form bounds}

Let $Q \in \mathbb{C}^{n\times n}$ be Hermitian positive--definite.
Let $\lambda_{\min}(Q)$ and $\lambda_{\max}(Q)$ denote its smallest
and largest eigenvalues. For any $x \in \mathbb{C}^n$,
\begin{equation}
  \lambda_{\min}(Q)\,\|x\|_2^2
  \le x^\ast Q x \le 
  \lambda_{\max}(Q)\,\|x\|_2^2.
  \label{eq:quadratic-bounds}
\end{equation}

\begin{lemma}[Resonance and norm equivalence]
\label{lem:resonance-norm}
Define $R(x) = -x^\ast Q x$ and $V(x) = x^\ast Q x$. Then
\begin{equation}
  V(x) = - R(x),
\end{equation}
and~\eqref{eq:quadratic-bounds} gives
\begin{equation}
  \lambda_{\min}(Q)\,\|x\|_2^2
  \le V(x) \le
  \lambda_{\max}(Q)\,\|x\|_2^2.
\end{equation}
\end{lemma}

\subsection{Drift calculation under Gaussian noise}

Assume for this subsection that $(\xi_t)$ is i.i.d.\ with
$\mathbb{E}[\xi_t] = 0$ and finite covariance
$\Sigma = \mathbb{E}[\xi_t\xi_t^\ast]$.

\begin{proposition}[One--step Lyapunov drift]
\label{prop:gaussian-drift}
Assume $T_t \equiv T$ is deterministic with $\|T\|_2 \le \kappa < 1$
and $\xi_t$ is i.i.d.\ with mean zero and covariance $\Sigma$. Let
$V(x) = x^\ast Q x$. Then
\begin{equation}
  \mathbb{E}[V(x_{t+1}) \mid x_t]
  = x_t^\ast T^\ast Q T x_t + \operatorname{tr}(Q\Sigma).
\end{equation}
In particular, if
\begin{equation}
  T^\ast Q T \preceq (1-\delta)\,Q
  \quad\text{for some } \delta \in (0,1),
  \label{eq:Q-contractive}
\end{equation}
then
\begin{equation}
  \mathbb{E}[V(x_{t+1}) \mid x_t]
  \le (1-\delta) V(x_t) + b,
  \quad b = \operatorname{tr}(Q\Sigma).
  \label{eq:gaussian-drift-ineq}
\end{equation}
\end{proposition}

\begin{proof}
Using $x_{t+1} = T x_t + \xi_t$,
\begin{align*}
  V(x_{t+1})
  &= (T x_t + \xi_t)^\ast Q (T x_t + \xi_t) \\
  &= x_t^\ast T^\ast Q T x_t
   + 2\,\mathrm{Re}\,(x_t^\ast T^\ast Q \xi_t)
   + \xi_t^\ast Q \xi_t.
\end{align*}
Taking conditional expectation given $x_t$ and using
$\mathbb{E}[\xi_t \mid x_t] = 0$, we obtain
\[
  \mathbb{E}[V(x_{t+1}) \mid x_t]
  = x_t^\ast T^\ast Q T x_t + \mathbb{E}[\xi_t^\ast Q \xi_t].
\]
Since $\xi_t$ has covariance $\Sigma$,
\[
  \mathbb{E}[\xi_t^\ast Q \xi_t]
  = \operatorname{tr}(Q\Sigma).
\]
If \eqref{eq:Q-contractive} holds,
\[
  x_t^\ast T^\ast Q T x_t
  \le (1-\delta)\, x_t^\ast Q x_t
  = (1-\delta) V(x_t),
\]
hence \eqref{eq:gaussian-drift-ineq}.
\end{proof}

\begin{remark}
Condition \eqref{eq:Q-contractive} is equivalent to requiring that
the operator norm of $Q^{1/2} T Q^{-1/2}$ is at most $\sqrt{1-\delta}$.
\end{remark}

\subsection{Drift with heavy--tailed noise}

When $\xi_t$ has Cauchy or more general $1$--stable components, the
second moment $\mathbb{E}[\xi_t^\ast Q \xi_t]$ is infinite and
Proposition~\ref{prop:gaussian-drift} does not apply. However, one can
still obtain a drift inequality on truncated Lyapunov functionals.

Fix a truncation level $L>0$ and define the truncated functional
\begin{equation}
  \tilde{V}_L(x)
  = \min\{V(x), L\}.
\end{equation}

\begin{proposition}[Truncated drift under Cauchy forcing]
\label{prop:cauchy-truncated}
Assume $T$ satisfies \eqref{eq:Q-contractive} and $\xi_t$ has symmetric
Cauchy components with scale $\gamma>0$. For each fixed $L>0$ there
exist constants $\delta_L \in (0,1)$ and $b_L < \infty$ such that
\begin{equation}
  \mathbb{E}[\tilde{V}_L(x_{t+1}) \mid x_t]
  \le (1-\delta_L)\, \tilde{V}_L(x_t) + b_L.
  \label{eq:truncated-drift}
\end{equation}
\end{proposition}

\begin{proof}[Proof sketch]
The proof follows a standard truncation argument. First note that for
bounded $x$ the distribution of $x_{t+1} = T x + \xi_t$ is a translated
Cauchy, which has heavy tails yet is locally integrable. For each
fixed $L$, $\tilde{V}_L$ is bounded and Lipschitz, hence the expectation
$\mathbb{E}[\tilde{V}_L(x_{t+1}) \mid x_t]$ is finite. One then splits
the analysis into two regions: $\{V(x_t) \le L/2\}$ and
$\{V(x_t) > L/2\}$, using the contraction of $T$ to dominate the
contribution of $x_t$ in the second region and the boundedness of
$\tilde{V}_L$ in the first region. The detailed constants are obtained
via a combination of the spectral bound on $T$ and crude bounds on
$\tilde{V}_L(Tx + \xi)$ for large $\|\xi\|$. The resulting inequality
has the form \eqref{eq:truncated-drift}.
\end{proof}

Proposition~\ref{prop:cauchy-truncated} is sufficient to deduce tightness
of $(x_t)$ and existence of stationary distributions under Cauchy noise
using standard Markov chain theory for truncation--based Lyapunov
functions.

\section{Distance and metric warping}

\subsection{Scalar metric factor}

Consider the baseline Riemannian structure $(L,g_0)$ with metric tensor
$g_0$ and Riemannian distance $d_0$. Let $m(\sigma)$ be a scalar metric
factor as in the main text, and define
\begin{equation}
  g_\sigma = m(\sigma)\, g_0.
\end{equation}
Then the corresponding distance $d_\sigma$ satisfies the following
simple scaling law.

\begin{lemma}[Distance scaling]
\label{lem:distance-scaling}
Let $\gamma : [0,1] \to L$ be a smooth curve with length
\[
\ell_0(\gamma) = \int_0^1 \sqrt{g_0(\dot{\gamma}(t),\dot{\gamma}(t))}\,dt
\]
under $g_0$ and length
\[
\ell_\sigma(\gamma) = \int_0^1 \sqrt{g_\sigma(\dot{\gamma}(t),\dot{\gamma}(t))}\,dt
\]
under $g_\sigma$. Then
\begin{equation}
  \ell_\sigma(\gamma) = \sqrt{m(\sigma)}\, \ell_0(\gamma),
\end{equation}
and hence
\begin{equation}
  d_\sigma(x,y) = \sqrt{m(\sigma)}\, d_0(x,y).
\end{equation}
\end{lemma}

\begin{proof}
For any tangent vector $v$,
\[
g_\sigma(v,v) = m(\sigma)\,g_0(v,v),
\]
so
\[
\sqrt{g_\sigma(\dot{\gamma}(t),\dot{\gamma}(t))}
 = \sqrt{m(\sigma)}\,
    \sqrt{g_0(\dot{\gamma}(t),\dot{\gamma}(t))}.
\]
Therefore
\[
\ell_\sigma(\gamma)
 = \int_0^1 \sqrt{m(\sigma)}\,
              \sqrt{g_0(\dot{\gamma}(t),\dot{\gamma}(t))}\,dt
 = \sqrt{m(\sigma)}\,\ell_0(\gamma).
\]
Taking an infimum over all curves $\gamma$ connecting $x,y$ yields the
distance scaling.
\end{proof}

\begin{remark}
In the main text we occasionally use $d_\sigma = d_0 / \sqrt{m(\sigma)}$
to refer informally to ``distance collapse''. The precise Riemannian
relation is as in \Cref{lem:distance-scaling}, i.e.\ distances scale
by $\sqrt{m(\sigma)}$. One can realize ``collapse'' by defining
$M(\sigma) = m(\sigma)^{-1} M_0$ instead; the choice is a modeling
convention and does not affect the qualitative analysis.
\end{remark}

\subsection{Warp thresholds}

Let $0 < \sigma_{\mathrm{loc}} < \sigma_{\mathrm{glob}}$ and set
\begin{equation}
  m(\sigma)
  =
  \begin{cases}
    1 + k_1 \sigma, & 0 \le \sigma < \sigma_{\mathrm{loc}},\\[0.2em]
    1 + k_1 \sigma_{\mathrm{loc}}
      + k_2(\sigma - \sigma_{\mathrm{loc}}),
      & \sigma_{\mathrm{loc}} \le \sigma < \sigma_{\mathrm{glob}},\\[0.2em]
    1 + k_1 \sigma_{\mathrm{loc}}
      + k_2(\sigma_{\mathrm{glob}} - \sigma_{\mathrm{loc}})
      + k_3(\sigma - \sigma_{\mathrm{glob}}),
      & \sigma \ge \sigma_{\mathrm{glob}},
  \end{cases}
\end{equation}
with $k_i > 0$. Then:

\begin{lemma}[Piecewise--linear curvature scaling]
\label{lem:piecewise-m}
The function $m$ is continuous, piecewise linear and strictly
increasing, with first--order curvature amplification at
$\sigma_{\mathrm{loc}}$ and $\sigma_{\mathrm{glob}}$. Consequently,
$d_\sigma$ is piecewise scaled relative to $d_0$ and the rate of
scaling increases when $\sigma$ crosses the thresholds.
\end{lemma}

\begin{proof}
Immediate from construction and \Cref{lem:distance-scaling}.
\end{proof}

This formalizes the notion of local vs global warp: for
$\sigma < \sigma_{\mathrm{loc}}$ the change in distance is mild;
for $\sigma_{\mathrm{loc}} \le \sigma < \sigma_{\mathrm{glob}}$ the
scaling rate increases (local dimpling); and for
$\sigma \ge \sigma_{\mathrm{glob}}$ the scaling rate increases again
(global warp).

\section{Tail behavior of stress under Gaussian vs Cauchy forcing}

Let $V(x) = x^\ast Q x$ and $\sigma_t = V(x_t)$ as above.

\subsection{Gaussian case}

Assume $\xi_t$ is Gaussian with mean zero and covariance $\Sigma$,
and that $T$ satisfies \eqref{eq:Q-contractive}. Under standard
irreducibility conditions, the Markov chain $(x_t)$ admits a unique
stationary distribution $\mu$ on $\mathbb{C}^n$ with finite second
moments.

\begin{proposition}[Exponential tails for $\sigma_t$ in Gaussian case]
\label{prop:gaussian-tails}
Under the above assumptions, there exists $\alpha>0$ and $C<\infty$
such that the stationary law of $\sigma_t$ satisfies
\begin{equation}
  \mathbb{P}_\mu(\sigma_t > u) \le C e^{-\alpha u}
\end{equation}
for all sufficiently large $u$.
\end{proposition}

\begin{proof}[Proof sketch]
Since the recursion is linear and the noise is Gaussian, the stationary
law $\mu$ is Gaussian with zero mean and covariance $S$ solving the
Lyapunov equation
\[
S = T S T^\ast + \Sigma.
\]
Thus $x_t \sim \mathcal{N}(0,S)$ under $\mu$. The quadratic form
$\sigma_t = x_t^\ast Q x_t$ is a generalized chi--square random
variable with a density that decays exponentially in $u$; see standard
results for quadratic forms in Gaussian vectors. The stated bound holds
for large $u$ with constants determined by the eigenvalues of $QS$.
\end{proof}

\subsection{Cauchy case}

Assume now that $\xi_t$ has symmetric Cauchy components with
scale $\gamma$. The recursion is an affine stochastic recurrence with
$1$--stable forcing. Exact characterizations of invariant laws are
nontrivial, but the following qualitative statement holds.

\begin{proposition}[Heavy tails for $\sigma_t$ in Cauchy case]
\label{prop:cauchy-tails}
Assume $\|T\|_2 < 1$ and $\xi_t$ has i.i.d.\ symmetric Cauchy
components with scale $\gamma>0$. Then any stationary distribution
$\mu$ of $(x_t)$ is heavy--tailed: for generic nonzero $u \in \mathbb{C}^n$
the marginal $u^\ast x_t$ has power--law tails of index $1$, and
$\sigma_t = x_t^\ast Q x_t$ has a tail which decays at most
polynomially in $u$.
\end{proposition}

\begin{proof}[Proof sketch]
For fixed nonzero $u$, the recursion for $y_t = u^\ast x_t$ is
\begin{equation*}
  y_{t+1} = u^\ast T x_t + u^\ast \xi_t.
\end{equation*}
Since linear images of Cauchy variables are Cauchy and the sum of
Cauchy variables is Cauchy, the marginal process $(y_t)$ is an affine
recursion with Cauchy forcing. In the stationary regime, $y_t$ is
Cauchy--distributed (up to a scale factor dependent on $T$ and $u$).
Cauchy tails decay as $|z|^{-1}$, hence the heavy--tail property.
The quadratic form $\sigma_t$ then inherits at most polynomial decay
from the coordinate projections; for large $u$,
$\mathbb{P}(\sigma_t > u)$ behaves like a power law.
\end{proof}

\begin{remark}
A complete description requires the theory of affine stochastic
recurrences with stable noise. The essential point for the present
framework is that the contraction $T$ prevents divergence of $x_t$,
but it does not suppress heavy tails introduced by the Cauchy forcing.
\end{remark}

\section{Summary of bounds}

For ease of reference, we summarize the main quantitative bounds
derived in this appendix:

\begin{itemize}
  \item Operator norm and spectral radius: $\rho(T_t) \le \|T_t\|_2
  \le \kappa < 1$, implying $\|T_{t-1}\cdots T_s\|_2 \le \kappa^{t-s}$.

  \item Trajectory bound:
  \[
  \|x_t\|_2 \le \kappa^t \|x_0\|_2 + 
  \sum_{j=0}^{t-1} \kappa^{t-1-j} \|\xi_j\|_2.
  \]

  \item Quadratic form equivalence:
  \[
  \lambda_{\min}(Q)\,\|x\|_2^2
  \le x^\ast Q x
  \le \lambda_{\max}(Q)\,\|x\|_2^2.
  \]

  \item Gaussian Lyapunov drift:
  \[
  \mathbb{E}[V(x_{t+1}) \mid x_t]
  \le (1-\delta)V(x_t) + \operatorname{tr}(Q\Sigma).
  \]

  \item Distance scaling under scalar metric:
  \[
  d_\sigma(x,y) = \sqrt{m(\sigma)}\, d_0(x,y).
  \]

  \item Gaussian case:
  exponential tails for $\sigma_t$ under stationarity.

  \item Cauchy case:
  polynomial (Cauchy--like) tails for linear functionals and for
  $\sigma_t$ under stationarity.
\end{itemize}

These estimates provide the analytic backbone for the NEL--$\Lambda_m$--Cauchy
framework developed in the main text.

\nocite{*}
\bibliographystyle{plain}
\bibliography{references}
\end{document}